\documentclass{leveldoc}
\DocUses{foundations}
\DocFingerprint{2ecc92e62c2d03fa}

\begin{document}

\section{What the level asks}\label{sec:asks}

The game's text:
\leveldescription

A positive particle is launched from the left edge of the board, straight to the right. The
detector sits in the upper right corner. The player places up to three fixed charges, each
of magnitude \num{1e6}, \num{2e6} or \num{4e6} and of either sign, on any nodes of the grid,
to bend the flight into the detector. The level is solved when the flight enters the detector
before it touches a charge, leaves the world or runs out of time (\xsecref{foundations}{sec:board}).

The reference solution places a single charge $Q = \val{Q}$ at the origin, the corner below
the launch point: it repels the particle upwards. This document derives the flight exactly:
the orbit of a relativistic charge about a fixed charge has a closed form.

\section{Setup}\label{sec:setup}

The level file, field by field (the appendix reproduces it):
\setuptable

In symbols (\xsecref{foundations}{sec:notation}): the particle has charge $q = \val{q}$, mass
$m = \val{m}$ and radius $a = 0$; the speed of light is $c = \val{c}$. It starts at
$\vect{x}_0 = (0, 10)$ along $+x$ with kinetic energy $T_0 = \val{Tzero}$. The reference charge
$Q = \val{Q}$ sits at $\vect{r}_\mathrm{c} = (0, 0)$ with radius $R = \val{R}$, and the coupling
is $\kappa = qQ = \val{kappa}$ (repulsive). The detector is the rectangle
$27 \le x \le 30$, $16 \le y \le 20$, whose left face lies at $x_D = \val{xD}$; the world is
$-1 \le x \le 31$, $-1 \le y \le 21$; the time limit is $t_{\max} = \val{tmax}$.

\begin{figure}[htbp]
\centering
\begin{tikzpicture}
\begin{axis}[axis equal image, width=0.92\linewidth, xmin=-1.6, xmax=31.6, ymin=-1.6, ymax=21.6,
  xtick={0,5,...,30}, ytick={0,5,...,20}, tick label style={font=\small},
  xlabel={$x$}, ylabel={$y$}, label style={font=\small}]
  \draw[black!8] (0,0) grid (30,20);
  \draw[black!45, dashed] (-1,-1) rectangle (31,21);
  \fill[green!50!black, fill opacity=0.18] (27,16) rectangle (30,20);
  \draw[green!40!black] (27,16) rectangle (30,20);
  \node[font=\small, green!30!black] at (28.5,18) {detector};
  \addplot[very thick, blue!70!black] table {orbit.dat};
  \fill[red!75!black] (0,0) circle[radius=0.3];
  \node[font=\small, red!60!black, anchor=south west] at (0.3,0.2) {$Q$};
  \draw[-{Stealth[length=6pt]}, thick] (0,10) -- (3,10);
  \fill (0,10) circle[radius=0.15];
  \node[font=\small, anchor=south] at (1.5,10.2) {launch};
  \fill[blue!70!black] (27,\rawval{yE}) circle[radius=0.2];
  \node[font=\small, anchor=east] at (26.6,\rawval{yE}) {$\vect{x}_\mathrm{e}$};
\end{axis}
\end{tikzpicture}
\caption{The board of the level with the reference solution: the charge $Q$ at the origin, the
exact orbit of \S\ref{sec:derivation} from the launch to the entry into the detector, and the
world's edge (dashed).}
\label{fig:board}
\end{figure}

\section{Principles and model}\label{sec:model}

\begin{principles}
  \item The field of a charged sphere outside it and the force on a non-polarizable particle,
    \xeqref{foundations}{eq:coulomb} and \xeqref{foundations}{eq:sphereforce}.
  \item The relativistic equation of motion \xeqref{foundations}{eq:eom} and the launch
    \xeqref{foundations}{eq:launch}.
  \item The conservation of energy \xeqref{foundations}{eq:energy} and of angular momentum in a
    central field \xeqref{foundations}{eq:angmom}; motion in the plane
    (\xsecref{foundations}{sec:plane}).
  \item The board's rules: the world, collisions, the detector, the outcome
    (\xsecref{foundations}{sec:board}).
\end{principles}

With one charge $Q$ at $\vect{r}_\mathrm{c}$, the particle's state $(\vect{x}, \vect{p})$ obeys
\begin{equation}\label{eq:eom-level}
  \deriv{\vect{x}}{t} = \frac{\vect{p}}{\gamma m}, \qquad
  \deriv{\vect{p}}{t} = \kappa\,\frac{\vect{x}-\vect{r}_\mathrm{c}}{\abs{\vect{x}-\vect{r}_\mathrm{c}}^3},
  \qquad \gamma = \sqrt{1 + \frac{\abs{\vect{p}}^2}{m^2c^2}},
\end{equation}
from $\vect{x}(0) = \vect{x}_0$, $\vect{p}(0) = \abs{\vect{p}_0}\,\uvec{x}$, until the first
event of \xsecref{foundations}{sec:board}. This is classical electrodynamics for this setup,
with these assumptions:

\begin{assumption}{The charge is held in place}\label{as:fixed}
The charge does not recoil: something holds it. This is a rule of the game, not an
approximation of the physics it describes.
\end{assumption}

\begin{assumption}{No radiation}\label{as:radiation}
The particle radiates while it is accelerated \xeqref{foundations}{eq:larmor}; the energy this
carries off is not taken from it. Along the reference flight the radiated energy is
$\val{Erad}$, a fraction $\val{EradFrac}$ of the launch energy (§\ref{sec:hand}, Step~7): it
changes the arrival point by a relative amount of the same order, far below anything that
matters here.
\end{assumption}

\begin{assumption}{A point particle}\label{as:point}
The particle's radius is $a = 0$, and it does not polarize: \xeqref{foundations}{eq:sphereforce}
is exact.
\end{assumption}

The motion stays in the plane $z = 0$ (\xsecref{foundations}{sec:plane}); from here on vectors
have two components.

\section{Derivation}\label{sec:derivation}

\subsection{Polar coordinates and the conserved quantities}

Take polar coordinates about the charge, $\vect{x} - \vect{r}_\mathrm{c} = r(\cos\varphi, \sin\varphi)$.
The velocity is $\dot r\,\uvec{r} + r\dot\varphi\,\uvec{\varphi}$, so the momentum has the
radial and azimuthal components $p_r = \gamma m\dot r$ and $p_\varphi = \gamma m r\dot\varphi$.
The angular momentum \xeqref{foundations}{eq:angmom} has only a $z$-component,
$L_z = r p_\varphi = \gamma m r^2\dot\varphi$, and it is constant. Unless the particle moves
along a radius ($L_z = 0$, §\ref{sec:others}), $\varphi$ therefore changes monotonically. Let
$\sigma = \operatorname{sign} L_z$, $L = \abs{L_z}$, and let
\begin{equation}\label{eq:theta}
  \theta = \sigma\,(\varphi - \varphi_0), \qquad \dot\theta = \frac{L}{\gamma m r^2} > 0,
\end{equation}
be the angle swept since the launch. The energy \xeqref{foundations}{eq:energy} with the
potential $\phi = Q/r$ is
\begin{equation}\label{eq:energy}
  W = \gamma mc^2 + \frac{\kappa}{r} = \gamma_0 mc^2 + \frac{\kappa}{r_0}.
\end{equation}

\subsection{The orbit equation}

Use $\theta$ as the variable along the orbit and $u = 1/r$, with $u' = \dd u/\dd\theta$. Then
$\dot r = (\dd r/\dd\theta)\,\dot\theta = -(u'/u^2)\,L u^2/(\gamma m)$, and
\begin{equation}\label{eq:radial}
  p_r = -L\,u', \qquad p_\varphi = \sigma L u.
\end{equation}
The energy relation $(\gamma mc^2)^2 = m^2c^4 + c^2(p_r^2 + p_\varphi^2)$ with
$\gamma mc^2 = W - \kappa u$ from \eqref{eq:energy} becomes
\begin{equation}\label{eq:energy-u}
  (W - \kappa u)^2 = m^2c^4 + c^2L^2\left(u'^2 + u^2\right).
\end{equation}
Differentiating with respect to $\theta$ gives
$-2\kappa(W - \kappa u)\,u' = 2c^2L^2\,u'(u'' + u)$; where $u' \ne 0$ we divide by
$2c^2L^2u'$, and since the solution is smooth the result holds where $u' = 0$ as well:
\begin{equation}\label{eq:orbit}
  u'' + \Gamma^2 u = -\frac{\kappa W}{c^2L^2}, \qquad \Gamma^2 = 1 - \frac{\kappa^2}{c^2L^2}.
\end{equation}
A linear equation with constant coefficients. For $\Gamma^2 > 0$ its solution from the launch
values $u_0 = 1/r_0$ and $u'_0 = -p_{r,0}/L$ \eqref{eq:radial} is
\begin{equation}\label{eq:general}
  u(\theta) = B + (u_0 - B)\cos\Gamma\theta + \frac{u'_0}{\Gamma}\sin\Gamma\theta,
  \qquad B = -\frac{\kappa W}{c^2L^2\Gamma^2},
\end{equation}
and for $\Gamma^2 < 0$ the same with $\cos\Gamma\theta \to \cosh\abs{\Gamma}\theta$ and
$(\sin\Gamma\theta)/\Gamma \to (\sinh\abs{\Gamma}\theta)/\abs{\Gamma}$. As $c\to\infty$, $\Gamma \to 1$ and
$W/c^2 \to m$: \eqref{eq:orbit} becomes Binet's equation $u'' + u = -\kappa m/L^2$ of the
Newtonian Kepler problem (\Goldstein{§3.5}). Relativity enters through $\Gamma < 1$: the orbit
turns more slowly in $\theta$ than a Newtonian one.

\subsection{The reference orbit}\label{sec:reference-orbit}

In the reference, $\vect{r}_\mathrm{c} = 0$ and $\vect{x}_0 = (0, 10)$: $r_0 = 10$ and
$\varphi_0 = \pi/2$. The launch momentum points along $+x$, perpendicular to the radius, so
$p_{r,0} = 0$: the launch is an apsis, $u'_0 = 0$, and $L_z = -r_0\abs{\vect{p}_0} < 0$: the
particle turns clockwise, $\sigma = -1$, $\theta = \pi/2 - \varphi$. Then \eqref{eq:general}
reads
\begin{equation}\label{eq:reference}
  u(\theta) = A\cos\Gamma\theta + B, \qquad A = u_0 - B,
\end{equation}
and the position is
\begin{equation}\label{eq:position}
  x(\theta) = \frac{\sin\theta}{u(\theta)}, \qquad y(\theta) = \frac{\cos\theta}{u(\theta)}.
\end{equation}
The energy relation \eqref{eq:energy-u} at the apsis ($u' = 0$, $u = A + B$) is a quadratic
for $A + B$; its positive root gives $A$ independently of the launch values,
\begin{equation}\label{eq:amplitude}
  A = \sqrt{B^2 + \frac{W^2 - m^2c^4}{c^2L^2\Gamma^2}},
\end{equation}
a check of the numbers (Step~3). Here $\kappa > 0$ makes $B < 0$, while $A + B = u_0 > 0$: $r$ grows
from $r_0$ without bound as $\theta$ approaches the asymptote
$\theta_\infty = \arccos(-B/A)/\Gamma$. The launch is the closest approach.

\subsection{Which event comes first}\label{sec:first-event}

By \eqref{eq:eom-level} with $\vect{r}_\mathrm{c} = 0$ and $\kappa > 0$,
$\dot p_x = \kappa x/r^3$ and $\dot p_y = \kappa y/r^3$. At the launch $x = 0$, $y = 10$,
$p_x = \abs{\vect{p}_0} > 0$, $p_y = 0$. While $x \ge 0$ and $y > 0$, neither component of
$\vect{p}$ decreases, so $p_x \ge \abs{\vect{p}_0} > 0$ and $p_y \ge 0$: $x$ increases and $y$
does not decrease, which keeps $x \ge 0$ and $y > 0$. So along the whole flight $x$ grows
strictly and $y$ never falls. Hence:
\begin{itemize}
  \item no collision: $r \ge r_0 = 10 > R + a$;
  \item the world's left and bottom edges are never reached; its right edge $x = 31$ only
    after $x = x_D$; its top edge $y = 21$ only after $y$ exceeds its value at $x = x_D$;
  \item the detector, which lies at $x \ge x_D$, can be entered first through its left face
    $x = x_D$, when $x$ reaches $x_D$, if $16 \le y \le 20$ there.
\end{itemize}
So the flight ends at the angle $\theta_\mathrm{e}$ where $x(\theta_\mathrm{e}) = x_D$, entering
the detector, provided $16 \le y(\theta_\mathrm{e}) \le 20$ and the time $t(\theta_\mathrm{e})$
is below $t_{\max}$ (Steps 4--6 show both). By \eqref{eq:position}, $\theta_\mathrm{e}$ solves
\begin{equation}\label{eq:crossing}
  f(\theta) = \sin\theta - x_D\left(A\cos\Gamma\theta + B\right) = 0,
\end{equation}
the condition $x = x_D$ multiplied by $u > 0$.

\subsection{The time}\label{sec:time}

From \eqref{eq:theta} and $\gamma m = (W - \kappa u)/c^2$,
\begin{equation}\label{eq:dtdth}
  \deriv{t}{\theta} = \frac{\gamma m r^2}{L} = \frac{W - \kappa u}{c^2 L u^2}.
\end{equation}
For the reference orbit \eqref{eq:reference}, with $s = \Gamma\theta$,
\begin{equation}\label{eq:time}
  t(\theta) = \frac{W\,I_2(\Gamma\theta) - \kappa\,I_1(\Gamma\theta)}{c^2 L\,\Gamma},
  \qquad I_n(s) = \int_0^s \frac{\dd s'}{(A\cos s' + B)^n}.
\end{equation}
Both integrals are elementary. With $\tau = \tan(s'/2)$,
$\cos s' = (1-\tau^2)/(1+\tau^2)$ and $\dd s' = 2\,\dd\tau/(1+\tau^2)$, so
$A\cos s' + B = \left((A+B) - (A-B)\tau^2\right)/(1+\tau^2)$ and, since $A > \abs{B}$ here,
\begin{equation}\label{eq:i1}
  I_1(s) = \int_0^{\tan(s/2)} \frac{2\,\dd\tau}{(A+B) - (A-B)\tau^2}
  = \frac{2}{\sqrt{A^2-B^2}}\,\operatorname{artanh}\!\left(\sqrt{\frac{A-B}{A+B}}\,\tan\frac{s}{2}\right),
\end{equation}
valid up to the asymptote, where the argument of $\operatorname{artanh}$ reaches $1$. And
because
$\dd/\dd s\left[\sin s/(A\cos s + B)\right] = (A + B\cos s)/(A\cos s + B)^2$ with
$A + B\cos s = (B/A)(A\cos s + B) + (A^2 - B^2)/A$,
\begin{equation}\label{eq:i2}
  I_2(s) = \frac{A\sin s}{(A^2-B^2)(A\cos s + B)} - \frac{B}{A^2-B^2}\,I_1(s).
\end{equation}

\subsection{The momentum at the entry}

At the entry, $\gamma_\mathrm{e} = (W - \kappa u_\mathrm{e})/(mc^2)$ by \eqref{eq:energy} and
$\abs{\vect{p}_\mathrm{e}} = mc\sqrt{\gamma_\mathrm{e}^2 - 1}$. Its direction follows from
\eqref{eq:radial} with $\uvec{r} = (\cos\varphi, \sin\varphi)$,
$\uvec{\varphi} = (-\sin\varphi, \cos\varphi)$:
\begin{equation}\label{eq:momentum}
  \vect{p} = p_r\,\uvec{r} + p_\varphi\,\uvec{\varphi}, \qquad p_r = -L\,u'(\theta), \quad
  p_\varphi = \sigma L\,u(\theta).
\end{equation}

\subsection{The radiated energy}

For \ref{as:radiation}: along the orbit the force is
$\vect{F} = \kappa(\vect{x}-\vect{r}_\mathrm{c})/r^3$, and from $\dot{\vect{p}} = \vect{F}$ with
$\dd(\gamma mc^2)/\dd t = \vect{v}\cdot\vect{F}$ the acceleration is
\begin{equation}\label{eq:acceleration}
  \vect{a} = \frac{1}{\gamma m}\left(\vect{F} - \frac{(\vect{v}\cdot\vect{F})\,\vect{v}}{c^2}\right).
\end{equation}
The power \xeqref{foundations}{eq:larmor} integrated over the flight,
$E_\mathrm{rad} = \int_0^{\theta_\mathrm{e}} P\,(\dd t/\dd\theta)\,\dd\theta$, is a smooth
integral, done with Simpson's rule \xeqref{foundations}{eq:simpson}.

\section{The reference solution by hand}\label{sec:hand}

\begin{procedure}
  \item\label{st:launch} \emph{The launch} \xeqref{foundations}{eq:launch}:
    $\gamma_0 = 1 + T_0/(mc^2) = \val{gammaZero}$,
    $\abs{\vect{p}_0} = \sqrt{T_0^2/c^2 + 2mT_0} = \val{pZero}$, and the speed
    $v_0 = \abs{\vect{p}_0}/(\gamma_0 m) = \val{vZero} = \val{betaZero}\,c$ (Newtonian
    mechanics would give $\sqrt{2T_0/m} = \val{vNewton}$).
  \item \emph{The conserved quantities} \eqref{eq:energy}: $r_0 = \val{rZero}$,
    $W = \gamma_0 mc^2 + \kappa/r_0 = \val{W}$, and $L = r_0\abs{\vect{p}_0} = \val{L}$.
  \item\label{st:constants} \emph{The orbit} \eqref{eq:orbit}, \eqref{eq:reference}:
    $\kappa/(cL) = \val{eta}$, so $\Gamma = \val{Gamma}$; $B = \val{B}$ and
    $A = u_0 - B = \val{A}$. The energy gives the same $A$ \eqref{eq:amplitude}, to
    $\val{Adiff}$. The asymptote lies at $\theta_\infty = \val{psiInf}$
    ($\val{psiInfDeg}^\circ$).
  \item\label{st:newton} \emph{The crossing of $x = x_D$}: Newton's method
    \xeqref{foundations}{eq:newton} on \eqref{eq:crossing}, with
    $f'(\theta) = \cos\theta + x_D A\Gamma\sin\Gamma\theta$, from $\theta = 1$:
    \begin{center}\small\val{newton}\end{center}
    The digits double each step: $\theta_\mathrm{e} = \val{psiE}$ ($\val{psiEdeg}^\circ$).
  \item \emph{The entry point} \eqref{eq:position}: $u_\mathrm{e} = A\cos\Gamma\theta_\mathrm{e} + B
    = \val{uE}$, $r_\mathrm{e} = \val{rE}$, and $y_\mathrm{e} = \cos\theta_\mathrm{e}/u_\mathrm{e}
    = \val{yE}$: between 16 and 20, so the particle enters the detector through its left face
    at $(\val{xD}, \val{yE})$ (§\ref{sec:first-event}).
  \item \emph{The time} \eqref{eq:time}--\eqref{eq:i2}: with $s_\mathrm{e} = \Gamma\theta_\mathrm{e}
    = \val{sE}$, $I_1 = \val{IOne}$ and $I_2 = \val{ITwo}$, so $t_\mathrm{e} = \val{tE}$,
    far within $t_{\max} = \val{tmax}$. (A quadrature of \eqref{eq:dtdth} gives the same to
    $\val{tEquadDiff}$.)
  \item\label{st:momentum} \emph{The momentum and the radiation}: $\gamma_\mathrm{e} = \val{gammaE}$
    (kinetic energy $\val{TE}$, gained from the repulsion), $\abs{\vect{p}_\mathrm{e}} = \val{pE}$,
    with $p_r = \val{prE}$ and $p_\varphi = \val{pphiE}$ \eqref{eq:momentum}:
    $\vect{p}_\mathrm{e} = (\val{pxE}, \val{pyE})$, $\val{angleE}^\circ$ above the $x$-axis.
    The radiated energy is $E_\mathrm{rad} = \val{Erad}$.
\end{procedure}

\begin{result}[The reference flight]
The particle enters the detector through its left face at $\vect{x}_\mathrm{e} = (\val{xD},
\val{yE})$ at $t_\mathrm{e} = \val{tE}$, with momentum
$\vect{p}_\mathrm{e} = (\val{pxE}, \val{pyE})$: the level is solved.
\end{result}

Every step is arithmetic with the functions $\sin$, $\cos$, $\tan$, $\arccos$ and
$\operatorname{artanh}$, which tables or series give to any accuracy; only Step~\ref{st:newton}
iterates, and its table shows how few iterations it needs. The radiated energy (Step
\ref{st:momentum}) needs a numerical integral; it matters only for \ref{as:radiation}.

\section{Comparison with the game}\label{sec:comparison}

\comparisontable

The document's numbers come from closed forms evaluated with 40 significant digits, so they are
exact for all purposes here. The game's differ from them by $\val{cmpTDiff}$ in the time and
$\val{cmpYDiff}$ in the entry height, the accuracy of its verified integration
(\xsecref{foundations}{sec:verification}), well inside the tolerance of $10^{-9}$ set from it.
The game's diagnostic of the radiated energy, its own quadrature of the same power, agrees to
$\val{cmpRadDiff}$ (relative). The table also checks three more placements in detail and every
single-charge placement's first event (§\ref{sec:others}).

\section{Other solutions}\label{sec:others}

\subsection{One charge anywhere}

With one charge of any value at any node, the flight is still the orbit about a single centre,
and \eqref{eq:general} with \eqref{eq:dtdth} gives it exactly; only the launch is no longer an
apsis ($u'_0 \ne 0$), and the orbit may be attractive ($\kappa < 0$). Three cases need care:
\begin{itemize}
  \item A charge on the launch line ($y = 10$, ahead of the launch) gives $L = 0$: the particle
    moves along the line. An attractive charge pulls it in (a collision); a repulsive one stops
    it at the distance $r_t$ where its kinetic energy vanishes, $W = mc^2 + \kappa/r_t$, i.e.
    \begin{equation}\label{eq:radial-motion}
      r_t = \frac{\kappa}{W - mc^2},
    \end{equation}
    and sends it back out through the world's left edge (unless $r_t \le R + a$).
  \item An attractive charge with $cL < \abs{\kappa}$ has $\Gamma^2 < 0$: $u$ grows like
    $\cosh\abs{\Gamma}\theta$ and the particle spirals into the charge, the relativistic
    \emph{fall to the centre} (\LandauII{§39}).
  \item An attractive charge with $W < mc^2$ binds the particle: the orbit is a rosette that
    precesses by $2\pi(1/\Gamma - 1)$ per turn, and runs until it touches the charge, leaves the
    world or enters the detector, or until $t_{\max}$.
\end{itemize}
To find the first event along an exact orbit, follow it in steps of
$\Delta\theta = \num{4e-4}$, test each chord against the detector, the charge's sphere of
radius $R + a$, the world and $t_{\max}$ (the time from \eqref{eq:dtdth} by the trapezoidal
rule), and take the first chord with an event. A chord departs from the arc by at most
$r\Delta\theta^2/8 \approx \num{1e-6}$ cells, so only a flight that passes that close to a
boundary could be misjudged; the agreement with the game on every placement (below) shows that
none does. Testing the chords, not only their end points, matters: placement~1 below enters the
detector $\val{cornerGap}$ cells from its corner and would leave it through the right face after
$\val{sliverAngle}$~rad, no more than $\val{sliverDepth}$ cells deep, a sliver that points
$\Delta\theta$ apart step over.

\begin{figure}[htbp]
\centering
\begin{tikzpicture}
\begin{groupplot}[group style={group size=3 by 2, horizontal sep=0.6cm, vertical sep=1.0cm},
  axis equal image, width=0.37\linewidth, xmin=-1, xmax=31, ymin=-1, ymax=21,
  xtick={0,10,20,30}, ytick={0,10,20}, tick label style={font=\scriptsize},
  title style={font=\small, yshift=-4pt}]
  \nextgroupplot[title={$Q = \num{1e6}$: \val{nPlusOne}}]
    \fill[green!50!black, fill opacity=0.2] (27,16) rectangle (30,20);
    \draw[black!40] (0,0) rectangle (30,20);
    \addplot[only marks, mark=*, mark size=1.1pt, red!70!black] table {map_plusone.dat};
    \fill (0,10) circle[radius=0.4];
  \nextgroupplot[title={$Q = \num{2e6}$: \val{nPlusTwo}}]
    \fill[green!50!black, fill opacity=0.2] (27,16) rectangle (30,20);
    \draw[black!40] (0,0) rectangle (30,20);
    \addplot[only marks, mark=*, mark size=1.1pt, red!70!black] table {map_plustwo.dat};
    \fill (0,10) circle[radius=0.4];
  \nextgroupplot[title={$Q = \num{4e6}$: \val{nPlusFour}}]
    \fill[green!50!black, fill opacity=0.2] (27,16) rectangle (30,20);
    \draw[black!40] (0,0) rectangle (30,20);
    \addplot[only marks, mark=*, mark size=1.1pt, red!70!black] table {map_plusfour.dat};
    \fill (0,10) circle[radius=0.4];
    \draw[blue!70!black, thick] (0,0) circle[radius=0.9];
  \nextgroupplot[title={$Q = \num{-1e6}$: \val{nMinusOne}}]
    \fill[green!50!black, fill opacity=0.2] (27,16) rectangle (30,20);
    \draw[black!40] (0,0) rectangle (30,20);
    \addplot[only marks, mark=*, mark size=1.1pt, blue!70!black] table {map_minusone.dat};
    \fill (0,10) circle[radius=0.4];
  \nextgroupplot[title={$Q = \num{-2e6}$: \val{nMinusTwo}}]
    \fill[green!50!black, fill opacity=0.2] (27,16) rectangle (30,20);
    \draw[black!40] (0,0) rectangle (30,20);
    \addplot[only marks, mark=*, mark size=1.1pt, blue!70!black] table {map_minustwo.dat};
    \fill (0,10) circle[radius=0.4];
  \nextgroupplot[title={$Q = \num{-4e6}$: \val{nMinusFour}}]
    \fill[green!50!black, fill opacity=0.2] (27,16) rectangle (30,20);
    \draw[black!40] (0,0) rectangle (30,20);
    \addplot[only marks, mark=*, mark size=1.1pt, blue!70!black] table {map_minusfour.dat};
    \fill (0,10) circle[radius=0.4];
\end{groupplot}
\end{tikzpicture}
\caption{Every node where a single charge solves the level, for each charge the player may
place (and how many nodes), from the exact orbits: repelling charges (red) below the launch
line, attracting ones (blue) above it. The black dot is the launch; the blue ring marks the
reference.}
\label{fig:map}
\end{figure}

Of the $\val{nNodes}$ nodes and six charge values, $\val{nSolving}$ single-charge placements
solve the level (Figure~\ref{fig:map}). Each family is a band: a charge repels the particle
upwards from below the launch line, or attracts it upwards from above, and the deflection must
fall within the angular window that leads into the detector; weaker charges must sit closer to
the flight. All $\val{cmpMapTotal}$ placements' first events agree with the game's verdicts
(Table~\ref{tab:comparison}, row \emph{every single-charge placement}).

Three of them in detail (\code{placements.json}), each by the exact orbit, its entry found by
Newton's method on the face it crosses:
\begin{enumerate}[label=\arabic*.]
  \item $Q = \val{detOneQ}$ at $\val{detOneNode}$: enters through the \val{detOneFace} face at
    $(\val{detOneX}, \val{detOneY})$, at $t = \val{detOneT}$, $\val{cornerGap}$ cells from the
    corner.
  \item $Q = \val{detTwoQ}$ at $\val{detTwoNode}$: through the \val{detTwoFace} face at
    $(\val{detTwoX}, \val{detTwoY})$, at $t = \val{detTwoT}$.
  \item $Q = \val{detThreeQ}$ at $\val{detThreeNode}$: through the \val{detThreeFace} face at
    $(\val{detThreeX}, \val{detThreeY})$, at $t = \val{detThreeT}$.
\end{enumerate}

\subsection{More charges}

The level allows three charges. With two or more the force is no longer central, the angular
momentum is not conserved, and the orbit has no closed form; the flight is then a numerical
solution of \eqref{eq:eom-level} with the fields added, which \levelref{slingshot} carries out.
Energy is still conserved \xeqref{foundations}{eq:energy}, which checks any such computation.

\section{What the game computes}\label{sec:engine}

The game integrates \eqref{eq:eom-level}, the sum of the placed charges' fields
\xeqref{foundations}{eq:coulomb} in the order of the level file, with its eighth-order
Runge--Kutta method at the tolerances $10^{-10}$ and $10^{-12}$ (\physref{5}); the reference
flight takes 19 steps in the verified run. It locates the entry into the detector on the
method's interpolant with a search that cannot miss a crossing (\physref{6}), and verifies
the outcome by comparing the two runs (\physref{7}). Its model is this document's: no
approximation beyond \ref{as:radiation}, whose size it reports from the same
power \xeqref{foundations}{eq:larmor} (\physref{8}). The engine's own tests check it against
this exact orbit: test T3 compares the deflection angle of relativistic Coulomb scattering,
$\theta = \abs{\pi - (2/\Gamma)\arccos(-B/A)}$, the asymptote of \eqref{eq:reference}, with the
integrated flights (\physref{9}).

\section{Further reading}\label{sec:reading}

\begin{itemize}
  \item \LandauII{§39}: the relativistic motion in a Coulomb field, the orbit
    \eqref{eq:general}, the precession of bound orbits and the fall to the centre.
  \item \Goldstein{§§3.5 and 3.10}: Binet's equation and Rutherford scattering, the Newtonian
    limit of this level.
  \item \Jackson{§12.1}: the relativistic equations of motion of a charge in external fields;
    §14.2 for the radiated power.
\end{itemize}

\levelfileappendix

\end{document}
